\documentclass[a0,portrait]{a0poster}

\usepackage{multicol} % This is so we can have multiple columns of text side-by-side
\columnsep=100pt % This is the amount of white space between the columns in the poster
\columnseprule=3pt % This is the thickness of the black line between the columns in the poster
\usepackage{fancybox}
\usepackage[svgnames]{xcolor} % Specify colors by their 'svgnames', for a full list of all colors available see here: http://www.latextemplates.com/svgnames-colors

\usepackage{pdfpages}
\usepackage{cite}
\usepackage{color}
\usepackage{colortbl}
\usepackage[cmex10]{amsmath}
\usepackage{amssymb}
\usepackage{soul}
\usepackage[normalem]{ulem}
\usepackage{booktabs}
\usepackage{color}
\usepackage{placeins}
\usepackage{upgreek}
\usepackage[utf8]{inputenc}
\usepackage{bm,bbm}
\usepackage{bbding}

\usepackage[export]{adjustbox}

%\usepackage{times} % Use the times font
%\usepackage{palatino} % Uncomment to use the Palatino font
\usepackage{graphicx} % Required for including images
\graphicspath{{../../../../Logos/}} % Location of the graphics files
\usepackage[font=small,labelfont=bf]{caption} % Required for specifying captions to tables and figures
\usepackage{amsfonts, amsmath, amsthm, amssymb} % For math fonts, symbols and environments
\usepackage{wrapfig} % Allows wrapping text around tables and figures
\usepackage[framemethod=TikZ]{mdframed}
\usepackage{titlesec} % Modify titles
%\usepackage{fontspec}
%\setmainfont{Calibri}
\newcommand{\ds}{\displaystyle}
\newcommand {\bo}[1]{\textbf{#1}}
\newcommand{\pa}[1]{\left({#1}\right)}
\newcommand{\co}[1]{\left[{#1}\right]}
\newcommand{\ch}[1]{\left\{{#1}\right\}}

\titleformat{\section}{\color{white}\normalfont\Large\bfseries}{\color{white}\thesection}{1em}{\colorbox{SteelBlue}}{}

\setlength{\columnseprule}{0pt}
\mdfdefinestyle{MyFrame}{%
	linecolor=SteelBlue,
	outerlinewidth=2pt,
	roundcorner=50pt,
	innertopmargin=\baselineskip,
	innerbottommargin=\baselineskip,
	innerrightmargin=20pt,
	innerleftmargin=20pt,
	backgroundcolor=white!50!white}
%\input{figconfig}
%\titleformat{command}[shape]{format}{label}{sep}{be %fore-code}{after-code}
\begin{document}
\begin{mdframed}[style=MyFrame]
%----------------------------------------------------------------------------------------
%	POSTER HEADER 
%----------------------------------------------------------------------------------------

% The header is divided into two boxes:
% The first is 75% wide and houses the title, subtitle, names, university/organization and contact information
% The second is 25% wide and houses a logo for your university/organization or a photo of you
% The widths of these boxes can be easily edited to accommodate your content as you see fit
\begin{minipage}[b]{0.33\linewidth}
\raggedright
\includegraphics[width=12cm,valign=t]{NZSA-logo-words-e1505785105112-1.png}
\end{minipage}
%
\begin{minipage}[b]{0.33\linewidth}
\centering
\hfill
\end{minipage}
% 
\begin{minipage}[b]{0.33\linewidth}
\raggedleft
\includegraphics[width=11cm,valign=t]{"Logo Offshore Standard Landscape Reversed RGB.png"}
\end{minipage}\\

\vspace{3cm}
\begin{minipage}[h]{0.98\linewidth}
\centering \huge \color{SteelBlue} \textbf{Ordinal Pattern Analysis for Early Bearing Fault Detection and Classification in Rotating Machinery -- First Results} \color{Black}\\ % Title
\Large \textbf{Nome S. Sobrenome\textsuperscript{1}, Nome S. Sobrenome\textsuperscript{2}, and Nome S. Sobrenome\textsuperscript{3}}\\ % Author(s)
\normalsize \textsuperscript{1,2} Universidade, Local, Brasil\textsuperscript{1,2} and Universidade, Local, Brasil\textsuperscript{3}\\ %[-0.5cm] % University/organization
email@email.com\textsuperscript{1}, email@email.br\textsuperscript{2} and email@email.com\textsuperscript{3}\\
\end{minipage}
\vspace{0.5cm} % A bit of extra whitespace between the header and poster content

%----------------------------------------------------------------------------------------

\begin{multicols}{2} % This is how many columns your poster will be broken into, a portrait poster is generally split into 2 columns


%----------------------------------------------------------------------------------------
%	ABSTRACT
%----------------------------------------------------------------------------------------

\section{Introduction}

Ordinal Patterns are transformations that encode the sorting characteristics of values in $\mathbbm R^D$ into $D!$ symbols.

They were proposed by Bandt \& Pompe in 2002, and have proven their adequacy at extracting valuable information about the system that produces the data.

One of the possible encodings is associated to the set of indexes that sort the $D$ values in non-decreasing order, for example:
$$
(4.2, 5.1, 7.1, 3.9, 8.6) \longmapsto (2, 3, 4, 1, 5) \longmapsto \otimes \text{ (one of }5!=120 \text{ symbols)}.
$$

A time series $\bm x = (x_1,x_2,\dots,x_{D+n-1})$ can be transformed into a sequence of symbols $\bm \pi=(\pi_1,\pi_2,\dots,\pi_n)$.
Then, we compute $\bm h =(h_1,h_2,\dots,h_{D!})$ the histogram of $\bm \pi$, and from it two descriptors:

\section{Data}\label{section:1}

\section{Methodology}\label{section2}

\section{Results}\label{section4}


\section{Conclusions and Future Work}

\section*{Acknowledgements}
%
Este trabalho foi parcialmente suportado por XXXXXX, com recursos  YYYYY.
%----------------------------------------------------------------------------------------
%	REFERENCES
%----------------------------------------------------------------------------------------

%\bibliographystyle{IEEEtran}
%\bibliography{IEEEabrv,references}

%---------------------------------------------------
-------------------------------------
\end{multicols}
\begin{center}
\color{SteelBlue}{New Zealand Statistical Association 2024 Conference}
\end{center}
\end{mdframed}
\end{document}